\chapter{总结与展望}

\section{全文工作总结}
本文围绕推荐系统中递归自改进（RSI）闭环的保真度控制问题展开：
（1）形式化了 RSI 闭环与两类失败模式，在 MovieLens-100K 上复现
了朴素自训练的崩塌并归因于头部自强化（命题~\ref{prop:collapse}）；
（2）提出双轴保真度门控——流行度分层置信标准化 + 合成分布 JS
预算——从单样本与集合分布两端控制自生成数据准入；（3）提出多视
图一致性伪标签生成，以扰动一致性替代模型 top-$k$ 偏好作为生成
依据；（4）在留一法全量排序协议下系统验证：完整框架将朴素自训
练 \ExpRounds{} 轮后的 \ExpNaiveEnd{} 提升至 \ExpOursEnd{}
（\ExpGainText{}），迭代曲线平稳无崩塌，消融证实各模块贡献独
立，全部实验种子固定、可由命令行逐字节复现。

\section{不足与局限}
客观陈述四点局限：（1）\textbf{增益上界}：RSI 增益来自骨干模型
未充分利用的预测信息，当骨干已接近其数据上限时框架退化为"无
害"而无法"增益"；（2）\textbf{生成开销}：一致性生成需要 $K$ 次
前向/用户，全量排序场景下约为训练开销的 $K$ 倍，工业规模需近似
或缓存；（3）\textbf{预算的刻画力}：分层 JS 距离约束的是一阶边
际分布，不能排除高阶（如物品共现结构）漂移；（4）\textbf{单数据
集验证}：跨域（电商/新闻/音乐）与更大规模（ML-1M、Amazon）的
稳健性未在本工作范围内检验。

\section{未来研究方向}
（1）\textbf{高阶保真度度量}：以物品共现图或序列 $n$-gram 分布
扩展漂移度量，覆盖结构级漂移；（2）\textbf{多模态扩展}：将一致
性生成扩展到多模态序列（文本评论、封面图），视图扰动天然存在于
模态缺失场景；（3）\textbf{在线部署}：预算机制与线上 A/B 反馈
的联动——漂移预算本身可作为线上安全阀；（4）\textbf{联邦场景}
：分布式客户端各自闭环、中心只聚合门控统计，隐私友好的自迭代
值得研究；（5）\textbf{可控与安全}：自生成数据的对抗鲁棒性
（恶意用户诱导闭环漂移）是落地前必须回答的问题。
